% BHU PYQ -- Manually transcribed & error-corrected
% Paper: BCH-311 : Advanced Company Accounts
% Exam: B.Com. (Hons.) Semester V Examination, 2022-23

\documentclass[11pt,a4paper]{article}
\usepackage[a4paper,top=2.5cm,bottom=2.5cm,left=2.8cm,right=2.8cm,headheight=14pt]{geometry}
\usepackage{amsmath,amssymb}
\usepackage[T1]{fontenc}
\usepackage[utf8]{inputenc}
\usepackage{lmodern,microtype,fancyhdr,enumitem,hyperref,lastpage,parskip}

\hypersetup{
    colorlinks=true,
    urlcolor=black,
    linkcolor=black,
    pdftitle={BCH-311: Advanced Company Accounts -- BHU B.Com. (Hons.) Sem V 2022-23}
}

\pagestyle{fancy}
\fancyhf{}
\renewcommand{\headrulewidth}{0.4pt}
\renewcommand{\footrulewidth}{0.4pt}
\fancyhead[L]{\small   \textit{Banaras Hindu University}}
\fancyhead[R]{\small   \textit{BCH-311: Advanced Company Accounts}}
\fancyfoot[C]{\small Page  \thepage\ of \pageref{LastPage}}


\newlist{parts}{enumerate}{2}
\setlist[parts,1]{label=(\alph*),leftmargin=*,itemsep=5pt,topsep=3pt}
\setlist[parts,2]{label=(\roman*),leftmargin=*,itemsep=3pt,topsep=2pt}

\newcommand{\pts}[1]{\hfill   \textit{[#1]}}


\begin{document}


\begin{center}
  {\large\bfseries BANARAS HINDU UNIVERSITY}\\[2pt]
  {
\normalsize B.Com. (Hons.) --- Semester V Examination, 2022-23}\\[6pt]
  \rule{0.8\linewidth}{0.5pt}\\[6pt]
  {\large\bfseries Paper No. BCH-311: Advanced Company Accounts}\\[4pt]
  {
\normalsize Time: 3 Hours\hfill Maximum Marks: 70}
\end{center}

\rule{\linewidth}{0.4pt}

\smallskip



\noindent   \textit{Instructions: Attempt all questions. All questions carry equal marks (14 marks each).}

\bigskip




\noindent   \textbf{Question 1.} \\
The extracts of various accounts from the books of X Ltd. as on 31 extsuperscript{st} March, 2021 are as follows:

\begin{center}
  
\begin{tabular}{|p{4.3cm}|r|p{4.3cm}|r|}
    \hline
   \textbf{Particulars} &   \textbf{Amount (Rs.)} &   \textbf{Particulars} &   \textbf{Amount (Rs.)} \\ \hline
    \raggedright 50,000 Equity Shares of Rs. 20 each & 10,00,000 & \raggedright Profit of the year & 10,00,000 \\
    \raggedright 10,000, 6\% Preference Shares of Rs. 100 each & 10,00,000 & \raggedright Loan from Bank & 50,000 \\
    \raggedright General Reserve & 8,00,000 & \raggedright S. Creditors & 2,00,000 \\
    \raggedright Stock & 6,80,000 & \raggedright Goodwill & 5,00,000 \\
    \raggedright Stores \& Spares & 1,30,000 & \raggedright Investments & 9,00,000 \\
    \raggedright S. Debtors & 6,76,000 & \raggedright Cash and Bank Balance & 1,80,000 \\
    \raggedright Land and Building & 2,80,000 & \raggedright Preliminary Expenses & 20,000 \\
    \raggedright Plant and Machinery & 6,60,000 & \raggedright Furniture \& Fittings & 24,000 \\ \hline
  \end{tabular}
\end{center}
X Ltd. is taken over by Y Ltd. under the following arrangements:

\begin{enumerate}[label=(\arabic*),leftmargin=*]
  \item All the assets and liabilities are taken over.
  \item Plant and Machinery is taken over at 10\% less.
  \item Furniture and Fittings are taken over at 20\% less.
  \item Investments and stock are taken over at 5\% higher value.
  \item X Ltd. settled an unprovided contingent liability of Rs. 30,000 at one-third of the value.
  \item The purchase consideration is to be met by issuing 4 fully paid equity shares of Rs. 10 each of Y Ltd. for each equity share in X Ltd., 20 fully paid equity shares of Rs. 10 each of Y Ltd. for each preference share in X Ltd., and Rs. 5,00,000 in cash.
\end{enumerate}
Show with full detail the calculation of purchase consideration and pass necessary entries in the books of X Ltd.\pts{14}

\centerline{   \textbf{OR}}




\noindent What is amalgamation in the nature of merger? Also give entries to record amalgamation in the nature of merger in the books of transferor company.\pts{14}

\medskip\rule{\linewidth}{0.2pt}




\noindent   \textbf{Question 2.} \\
The balances of various accounts of Roogna Ltd. as on 31 extsuperscript{st} March, 2022 are as follows:

\begin{center}
  
\begin{tabular}{|p{4.3cm}|r|p{4.3cm}|r|}
    \hline
   \textbf{Accounts with Credit Balances} &   \textbf{Amount (Rs.)} &   \textbf{Accounts with Debit Balances} &   \textbf{Amount (Rs.)} \\ \hline
    \raggedright 10,000 Equity Shares of Rs. 100 each & 10,00,000 & \raggedright Goodwill & 2,25,000 \\
    \raggedright 5,000, 10\% Preference Shares of Rs. 100 each & 5,00,000 & \raggedright Land and Building & 4,50,000 \\
    \raggedright 9\% Debentures & 8,00,000 & \raggedright Plant and Machinery & 4,05,000 \\
    \raggedright Sundry Creditors & 3,00,000 & \raggedright Patents & 3,00,000 \\
    & & \raggedright Stock & 2,60,000 \\
    & & \raggedright Debtors & 3,10,000 \\
    & & \raggedright Cash in Hand & 1,50,000 \\
    & & \raggedright Preliminary Expenses & 2,00,000 \\
    & & \raggedright Profits and Loss A/C & 3,00,000 \\ \hline
   \textbf{Total} &   \textbf{26,00,000} &   \textbf{Total} &   \textbf{26,00,000} \\ \hline
  \end{tabular}
\end{center}
The following scheme of internal reconstruction is approved:

\begin{parts}
  \item Equity shares are to be reduced to shares of Rs. 50 each as fully paid.
  \item Preference shares are to be reduced by 40\% and the rate of dividend is to be increased to 12\%.
  \item The value of land and building is to be increased by 20\%.
  \item All fictitious assets are to be eliminated. A Patent right that was acquired for Rs. 50,000 is now of no use.
  \item Creditors have agreed to waive 20\% of their total claim if they are paid 50\% of the balance amount immediately.
\end{parts}
Pass necessary entries to record reconstruction of the company. Apply your reason regarding goodwill.\pts{14}

\centerline{   \textbf{OR}}




\noindent Differentiate between alteration of share capital and reduction of share capital.\pts{14}

\medskip\rule{\linewidth}{0.2pt}




\noindent   \textbf{Question 3.} \\
Explain the calculation of cost of control and minority interest in the preparation of consolidated balance sheet of Holding Company.\pts{14}

\centerline{   \textbf{OR}}




\noindent From the following Balance Sheets of H. Ltd. and S. Ltd., prepare the consolidated Balance Sheet as at 31-03-2022:

\begin{center}
  
\begin{tabular}{|p{6.0cm}|r|r|}
    \hline
   \textbf{Particulars} &   \textbf{H. Ltd. (Rs.)} &   \textbf{S. Ltd. (Rs.)} \\ \hline
   \textbf{I. Equity and Liabilities} & & \\
   \textbf{(1) Shareholder's Funds} & & \\
    \quad Equity Share Capital (Shares of Rs. 100 each) & 10,00,000 & 4,00,000 \\
    \quad   \textbf{Reserves and Surplus:} & & \\
    \quad P \& L A/C & 6,00,000 & 3,00,000 \\
    \quad General reserve & 1,00,000 & -- \\
   \textbf{(2) Non-Current Liabilities} & 5,00,000 & 1,00,000 \\
   \textbf{(3) Current Liabilities} & & \\
    \quad Sundry Creditors & 3,00,000 & 50,000 \\ \hline
   \textbf{Total} &   \textbf{25,00,000} &   \textbf{8,50,000} \\ \hline
   \textbf{II. Assets} & & \\
   \textbf{(1) Non-current Assets} & & \\
    \quad Sundry Assets & 14,50,000 & 7,00,000 \\
    \quad Investments (3,000 Shares of S. Ltd.) & 7,00,000 & -- \\
   \textbf{(2) Current Assets} & & \\
    \quad Cash and Other Current Assets & 3,50,000 & 1,50,000 \\ \hline
   \textbf{Total} &   \textbf{25,00,000} &   \textbf{8,50,000} \\ \hline
  \end{tabular}
\end{center}
The shares in S. Ltd. were acquired by the H. Ltd. on 1-4-2021. On that date, the P \& L A/C of the S. Ltd. showed a balance of Rs. 50,000.\pts{14}

\medskip\rule{\linewidth}{0.2pt}




\noindent   \textbf{Question 4.} \\
A winding up order was issued against Alfa Ltd. and the following information is obtained with regard to the assets and liabilities as on 30 extsuperscript{th} June, 2006:

\begin{itemize}[leftmargin=*]
  \item Land and Building (book value Rs. 4,50,000) valued at Rs. 3,75,000.
  \item First Mortgage of Land and Building: Rs. 3,00,000.
  \item Second Mortgage of Land and Building: Rs. 1,12,500.
  \item Managing Director's salary unpaid (6 months): Rs. 22,500.
  \item Staff Salary Unpaid (one month): Rs. 16,050.
  \item Debtors (Good): Rs. 31,500.
  \item Debtors (Bad): Rs. 72,750.
  \item Doubtful Debtors (estimated to realise 50 percent): Rs. 12,900.
  \item Bank Overdraft (unsecured): Rs. 58,125.
  \item Cash in Hand: Rs. 825.
  \item Plant and Machinery (book value Rs. 2,47,500), estimated to realize Rs. 1,74,000.
  \item Stock (at cost Rs. 50,850), estimated to realize Rs. 33,900.
  \item Issued Capital: Equity shares of Rs. 10 each fully called-up: Rs. 1,50,000.
  \item Calls in arrears: Rs. 3,000 (estimated to realise Rs. 1,500).
  \item Unsecured creditors: Rs. 2,96,250.
  \item 8\% Debentures carrying a floating charge on the undertaking: Rs. 1,50,000.
  \item Interest on debentures for 1 extsuperscript{st} September and 1 extsuperscript{st} April were paid on due dates.
  \item A Contingent liability for damages of Rs. 37,500 is estimated to be settled for Rs. 18,000.
  \item Income-tax liability: For 30 June, 2004: Rs. 5,250; For 30 June, 2005: Rs. 1,275; For 30 June, 2006: Rs. 2,700.
  \item The reserves of the company on 1 extsuperscript{st} July, 2005 amounted to Rs. 7,500.
\end{itemize}
You are required to prepare a statement of affairs.\pts{14}

\centerline{   \textbf{OR}}




\noindent Write a note on any one of the following:\pts{14}

\begin{parts}
  \item Deficiency Account
  \item Liquidator's Final Statement of Account
\end{parts}

\medskip\rule{\linewidth}{0.2pt}




\noindent   \textbf{Question 5.} \\
What is ESOP? Give the accounting entries to record ESOP when it is used as a share based payment system.\pts{14}

\centerline{   \textbf{OR}}




\noindent What is earning per share? How is it calculated?\pts{14}

\vfill
\rule{\linewidth}{0.4pt}

\begin{center}\small
\href{https://bhu-exam.vercel.app/}{Click here to visit our website}\\[4pt]
\href{https://me-aryan.vercel.app/}{Click here to visit our contributor}\\[4pt]
\href{https://quashan-me.vercel.app/}{Click here to visit our contributor}
\end{center}

\end{document}
